\rok{Новембар 2023.}{B}

Први практични тест из Алгоритама и структура података

\zadatak{1}{Bubble Sort}
Написати \textit{Bubble Sort} алгоритам за сортирање целобројног низа дужине $n$.

\textbf{Додатне напомене за решавање задатка:}
\begin{itemize}
    \item Улаз је несортиран низ целих бројева.
    \item Излаз је сортирани низ целих бројева.
    \item У \textit{template}-у се налази класа \texttt{RandomArrayGenerator} коју треба искористити за генерисање низа случајних бројева.
    \item Креирати класу \texttt{BubbleSort} и у оквиру ње генерисати методу:
    \begin{Verbatim}[fontsize=\footnotesize, numbers=left, xleftmargin=10mm]
public static void sort(long[] arr)
    \end{Verbatim}
\end{itemize}



\zadatak{2}{Инвертовање једноструко повезане листе}
Креирање функције која ће инвертовати једноструко повезану листу. Алгоритам је неопходно написати са линеарном временском комплексношћу $O(n)$.

\textbf{Додатни захтеви за решавање задатка:}
\begin{itemize}
    \item Алгоритам писати у оквиру класе \texttt{LinkedList}.
    \item Захтеван потпис методе:
    \begin{Verbatim}[fontsize=\footnotesize, numbers=left, xleftmargin=10mm]
public void InvertList()
    \end{Verbatim}
\end{itemize}

\textbf{Примери:}
\begin{itemize}
    \item \textbf{Улаз:} $2 \rightarrow 3 \rightarrow 1 \rightarrow 6 \rightarrow 4$ \\
          \textbf{Излаз:} $4 \rightarrow 6 \rightarrow 1 \rightarrow 3 \rightarrow 2$
    \item \textbf{Улаз:} $2 \rightarrow 3$ \\
          \textbf{Излаз:} $3 \rightarrow 2$
\end{itemize}



\zadatak{3}{Брисање узастопних парова из низа бројева}
Креирати алгоритам који из низа целих бројева брише све узастопне парове вредности. Узастопни парови подразумевају вредности које се налазе на суседним позицијама у низу. Елементи из низа се бришу искључиво уколико се у проласку кроз низ два иста броја нађу један поред другог.

\textbf{Додатни захтеви за решавање задатка:}
\begin{itemize}
    \item Задатак решавати помоћу структуре стека (класа \texttt{Stack}) која је дата у оквиру шаблона задатка.
    \item Алгоритам треба да резултује \texttt{bool} вредношћу, у зависности од тога да ли су обрисани сви узастопни парови.
    \item Захтеван потпис методе:
    \begin{Verbatim}[fontsize=\footnotesize, numbers=left, xleftmargin=10mm]
public bool ShortenArray(int[] array)
    \end{Verbatim}
\end{itemize}

\textbf{Пример 1:}
\begin{itemize}
    \item \textbf{Улаз:} \texttt{[1, 2, 2, 3, 3, 1]}
    \item \textbf{Кораци брисања елемената у низу:}
    \begin{itemize}
        \item \texttt{[1\ 2\ 2\ 3\ 3\ 1]}
        \item \texttt{[2\ 2\ 3\ 3]}
        \item \texttt{[3\ 3]}
        \item \texttt{[]}
    \end{itemize}
    \item \textbf{Излаз:} \texttt{true}
\end{itemize}

\textbf{Пример 2:}
\begin{itemize}
    \item \textbf{Улаз:} \texttt{[2, 2, 1, 4, 3, 1]}
    \item \textbf{Кораци брисања елемената у низу:}
    \begin{itemize}
        \item \texttt{[2\ 2\ 1\ 4\ 3\ 1]}
        \item \texttt{[1\ 4\ 3\ 1]}
    \end{itemize}
    \item \textbf{Излаз:} \texttt{false}
\end{itemize}

\sablon{AISP2023_Test1_GrupaB}{2023/Novembar/GrupaB/AISP2023_Test1_GrupaB.linq}
